\section{Face splitting}\label{sec:split}
Think of $C$ as the part of a source whose outcomes are seen, and $D$ as the part that sees them.
If $D$ can select its steered state from a catalogue, the two parts can be supplied separately.

\begin{rulebox}{Face splitting for a closed family of sources}
\begin{theorem}\label{thm:split}
Let $D$ be a non-empty set of nodes of $\G$ and let $F_1,\dots,F_k$ $(k\ge 1)$ be distinct sources that strictly
contain $D$. Put $C_i=F_i\setminus D$ and $C=C_1\cup\dots\cup C_k$; the other sources of $\G$ are the
\emph{remaining} ones. Assume
\begin{description}[leftmargin=1.9cm,style=nextline,itemsep=1pt]
\item[(S1) closure] no remaining source contains a node of $C$ (the family $F_1,\dots,F_k$ is \emph{closed});
\item[(S2) visibility] every node of $D$ sees all nodes of $C$ and all their settings:
  $C\cup\Pa(C)\subseteq\Pa(d)$ for every $d\in D$, where $\Pa(C)=\bigcup_{c\in C}\Pa(c)$.
\end{description}
Let $\G'$ be obtained from $\G$ by replacing the sources $F_1,\dots,F_k$ by $C_1,\dots,C_k$ and $D$, and then
discarding every set that has a single member or is contained in another source. Then
\[
\Q(\G')=\Q(\G)\qquad\text{and}\qquad\C(\G')=\C(\G).
\]
\end{theorem}
\end{rulebox}

\begin{figure}[H]
\centering
\begin{tikzpicture}[mdag]
  \begin{scope}
    \vn{y}{-1.5,1.0}{$y$}
    \vn[Cn]{c1}{0,2}{$c_1$} \vn[Cn]{c2}{2,2}{$c_2$} \vn[Cn]{c3}{4,2}{$c_3$} \vn[Dn]{d}{2,0}{$d$}
    \qs[splitsrc]{f1}{1.0,1.5} \qs[splitsrc]{f2}{3.0,1.5}
    \feed[splitedge]{f1}{c1,c2,d} \feed[splitedge]{f2}{c2,c3,d}
    \de{c1}{d} \de{c2}{d} \de{c3}{d} \de{y}{c1} \de{y}{d}
    \node[lab,text=splitcol!80!black] at (0.45,1.05) {$F_1$}; \node[lab,text=splitcol!80!black] at (3.55,1.05) {$F_2$};
  \end{scope}
  \node[rel] at (5.4,1.0) {\equivsign};
  \begin{scope}[shift={(8.2,0)}]
    \vn{y}{-1.5,1.0}{$y$}
    \vn[Cn]{c1}{0,2}{$c_1$} \vn[Cn]{c2}{2,2}{$c_2$} \vn[Cn]{c3}{4,2}{$c_3$} \vn[Dn]{d}{2,0}{$d$}
    \qs{f1}{1.0,2.75} \qs{f2}{3.0,2.75}
    \feed{f1}{c1,c2} \feed{f2}{c2,c3}
    \de{c1}{d} \de{c2}{d} \de{c3}{d} \de{y}{c1} \de{y}{d}
    \node[lab,text=qcol] at (0.45,2.95) {$C_1$}; \node[lab,text=qcol] at (3.55,2.95) {$C_2$};
  \end{scope}
\end{tikzpicture}
\caption{Two overlapping sources are split along $D=\{d\}$. The green receiver sees $c_1,c_2,c_3$ and the setting $y$. Splitting both orange sources leaves the blue nodes sharing the two smaller sources.}
\label{fig:split-generic}
\end{figure}

\begin{figure}[H]
\centering
\begin{tikzpicture}[mdag]
  \begin{scope}
    \vn[Cn]{a}{0,0}{$a$} \vn[Dn]{b}{2.4,0}{$b$} \qs[splitsrc]{q}{1.2,1.0}
    \feed[splitedge]{q}{a,b} \de{a}{b}
    \node[tag] at (1.2,-0.7) {a source on an edge};
  \end{scope}
  \node[rel] at (3.5,0.2) {$=$};
  \begin{scope}[shift={(4.6,0)}]
    \vn[Cn]{a}{0,0}{$a$} \vn[Dn]{b}{3.6,0}{$b$}
    \de{a}{b} \node[lab,above] at (1.1,0) {$x_a$};
    \node[draw=qcol, fill=qcol!12, rounded corners=1.5pt, minimum width=8mm, minimum height=4.2mm, font=\scriptsize] (r) at (0,1.05) {$\rho_a$};
    \draw[qedge] (r) -- (a);
    \foreach \i/\lbl in {0/{$\tau_{0}$},1/{$\tau_{1}$},2/{$\tau_{2}$}} {
      \node[draw=qcol, fill=qcol!12, rounded corners=1.5pt, minimum width=8mm, minimum height=4.2mm, font=\scriptsize]
        (e\i) at ({2.75+0.85*\i},1.05) {\lbl};
      \draw[qedge] (e\i) -- (b);
    }
    \node[tag] at (1.8,-0.7) {private states: $a$'s reduced state, $b$'s catalogue};
  \end{scope}
  \node[rel] at (9.6,0.2) {$=$};
  \begin{scope}[shift={(10.7,0)}]
    \vn[Cn]{a}{0,0}{$a$} \vn[Dn]{b}{2.4,0}{$b$} \de{a}{b}
    \node[tag] at (1.2,-0.7) {no source};
  \end{scope}
\end{tikzpicture}
\caption{The smallest case of \cref{thm:split}: $C=\{a\}$, $D=\{b\}$, $k=1$. The entry $\tau_x$ of the catalogue is
the state that the outcome $x$ of $a$ would have left on the share of $b$.}
\label{fig:split-idea}
\end{figure}

With just $a\to b$ and the source $\{a,b\}$, measuring $a$ leaves $b$ a state determined by $x_a$.
The edge tells $b$ which state to prepare locally. For several receivers this private preparation becomes
a shared catalogue. If $a$ also has a setting, every receiver needs that setting to select its entry.

\begin{proof}[Proof of \cref{thm:split}]
Structural dominance gives one inclusion. For the other, fix a model of the unsplit family and put
$Y=\Pa(C)\setminus C$, $z=(x_C,x_Y)$. By (S1) the product $M_C(z)$ of the measurements at $C$
acts only on shares of the $F_i$. After tracing out those shares, the operator left on $D$ is
\[
 \tau_z=\Tr_C\!\left[(M_C(z)\otimes\id)\bigotimes_i\rho_{F_i}\right],
 \qquad q_z=\Tr\tau_z.
\]
The trace $q_z$ is also the weight of $x_C$ at those settings when each $C_i$ receives its own reduced
state $\Tr_D\rho_{F_i}$. Supply these reduced states independently and keep the old measurements at $C$.
Let the new source on $D$ emit a catalogue with entry $\tau_z/q_z$ whenever $q_z>0$, and a fixed state
otherwise. By (S2) each $d\in D$ reads $z$ and measures the matching entry together with its other shares.
A singleton can prepare its entry privately, and a source containing $D$ can distribute the catalogue.

For each assignment in \cref{eq:trace}, the weight $q_z$ times the selected state gives exactly $\tau_z$.
Zero weight contributes zero in both models. Everything outside the family sees the same visible values
and retains its old shares. Diagonal inputs give diagonal entries, proving the classical equality too.
\end{proof}

\begin{validbox}[label=ex:setting]{a setting that both nodes see}
\begin{minipage}{0.55\linewidth}
$D=\{d\}$, $F_1=\{c,d\}$, $C=\{c\}$, and $c$ has a setting $y$. (S2) asks that $d$ see $y$ as well, which
holds:\\[3pt]
\str{y→c, y→d, c→d; Q\{c,d\}} \ $\equiv$ \ \str{y→c, y→d, c→d}.\\[3pt]
The catalogue is indexed by $(x_c,x_y)$. The edge $y\to d$ supplies the second index.
\end{minipage}\hfill
\begin{minipage}{0.42\linewidth}\centering
\begin{tikzpicture}[mdag]
  \vn{y}{0,0.65}{$y$} \vn[Cn]{c}{1.3,1.3}{$c$} \vn[Dn]{d}{1.3,0}{$d$} \qs[splitsrc]{q}{2.1,0.65}
  \feed[splitedge]{q}{c,d} \de{y}{c} \de{y}{d} \de{c}{d}
  \node[rel] at (2.95,0.65) {\equivsign};
  \begin{scope}[shift={(3.9,0)}]
    \vn{y}{0,0.65}{$y$} \vn[Cn]{c}{1.3,1.3}{$c$} \vn[Dn]{d}{1.3,0}{$d$}
    \de{y}{c} \de{y}{d} \de{c}{d}
  \end{scope}
\end{tikzpicture}
\end{minipage}
\end{validbox}

\begin{validbox}[label=ex:together]{two sources that must be split together}
\begin{minipage}{0.5\linewidth}
$D=\{d\}$, $F_1=\{a,b,d\}$, $F_2=\{a,c,d\}$, so $C_1=\{a,b\}$ and $C_2=\{a,c\}$ have the node $a$ in common, which
is also a parent of $b$ and $c$.\\[3pt]
\str{a→b, a→c, a→d, b→d, c→d; Q\{a,b,d\}, Q\{a,c,d\}} \ $\equiv$\\
\str{a→b, a→c, a→d, b→d, c→d; Q\{a,b\}, Q\{a,c\}}.\\[3pt]
Here $a$ may measure its two shares jointly. Once $a$, $b$ and $c$ have all measured, their outcomes determine
an entry of the catalogue on the two shares of $d$. Neither source satisfies (S1) on its own, and
with $k=1$ alone the two structures stay in different classes.
\end{minipage}\hfill
\begin{minipage}{0.47\linewidth}\centering
\begin{tikzpicture}[mdag,scale=0.76]
  \newcommand{\togethernodes}{\vn[Cn]{a}{1.5,2.3}{$a$} \vn[Cn]{b}{-0.2,1.1}{$b$} \vn[Cn]{c}{3.2,1.1}{$c$}
    \vn[Dn]{d}{1.5,0}{$d$} \de{a}{d} \de{b}{d} \de{c}{d}
    \draw[edge] (a) to[bend right=32] (b); \draw[edge] (a) to[bend left=32] (c);}
  \togethernodes
  \qs[splitsrc]{f1}{0.95,1.15} \qs[splitsrc]{f2}{2.05,1.15}
  \feed[splitedge]{f1}{a,b,d} \feed[splitedge]{f2}{a,c,d}
  \node[rel] at (4.15,1.1) {\equivsign};
  \begin{scope}[shift={(5.3,0)}]
    \togethernodes
    \qs{f1}{0.95,1.15} \qs{f2}{2.05,1.15}
    \feed{f1}{a,b} \feed{f2}{a,c}
  \end{scope}
\end{tikzpicture}
\end{minipage}
\end{validbox}

\paragraph{Comparison with classical face splitting.}
The classical rules of Ansanelli, Wolfe and Spekkens \cite{AWS} share the visibility condition (S2).
They differ in which other sources may meet $C$. The last column gives the number of four-node classes
if that splitting rule is used alongside the quantum edge rules.

\begin{center}\small
\begin{tabular}{@{}lp{8cm}cr@{}}
\toprule
Splitting rule & Other sources allowed to meet $C$ & Quantum valid & Classes\\
\midrule
Weak & None for each individual split & yes & 461\\
Closed family & Other sources split along with it & yes & 446\\
Moderate & Remaining sources that contain $D$ & no & 450\\
Strong & Both kinds & no & 438\\
\bottomrule
\end{tabular}
\end{center}
The joint split above goes beyond the weak rule. The six-node witness in \cref{sec:leaky} rules out
moderate splitting in general, including as a special case of the strong rule. For a fixed $D$ the
code need only split the largest closed family satisfying (S2). Splitting fewer sources gives an
intermediate structure covered by the sandwich (\cref{cor:sandwich}).
